% Warmup: worked example of the base witness pair.

\section{A worked example: two triangles versus the six-cycle}\label{sec:warmup}

Before stating the dichotomy in general, this chapter works out one
pair of graphs in complete detail. The pair will be used as a witness
in every non-membership claim in Section~\ref{sec:measures} and as the
building block for the distributional witness of
Section~\ref{sec:entanglement}.

\subsection{The two graphs}

Fix the vertex set $\vertices = \{1, 2, 3, 4, 5, 6\}$. Let $A$ be the
disjoint union of two triangles,
\[
A = C_3 \sqcup C_3, \qquad
\edges(A) = \{\{1,2\},\ \{1,3\},\ \{2,3\},\ \{4,5\},\ \{4,6\},\ \{5,6\}\}.
\]
Let $B$ be the six-cycle,
\[
B = C_6, \qquad
\edges(B) = \{\{1,2\},\ \{2,3\},\ \{3,4\},\ \{4,5\},\ \{5,6\},\ \{1,6\}\}.
\]
Both graphs have six vertices and six edges. Every vertex in either
graph has degree two. The two graphs are not isomorphic: $A$ has
two connected components while $B$ has one.

\subsection{1-WL cannot distinguish $A$ from $B$}

Assign the constant feature $x_v \equiv \bot$ to every vertex of
either graph. Colour refinement proceeds as follows.

\emph{Round zero.} Every vertex receives the initial colour
$c_v^{(0)} = \tau(\bot) = \alpha$ on both $A$ and $B$.

\emph{Round one.} A vertex $v$ of $A$ or $B$ has two neighbours, both
coloured $\alpha$, so
\[
c_v^{(1)} = \tau(\alpha, \{\!\{\alpha, \alpha\}\!\}) = \beta.
\]
The colour $\beta$ is assigned to every vertex of either graph.

\emph{Higher rounds.} Round two proceeds identically to round one with
$\beta$ in place of $\alpha$. Refinement has stabilised.

The graph-level fingerprints are
\[
\mathrm{hist}(A, \bot) = \{\!\{\beta,\beta,\beta,\beta,\beta,\beta\}\!\}
= \mathrm{hist}(B, \bot),
\]
so $A \sim_{\wl} B$. Colour refinement cannot separate the two
graphs.

By the upper-bound direction of Theorem~\ref{thm:dichotomy}, no MPNN with
continuous update and readout can separate them either. For every $M
\in \mpnn$, $M(A, \bot) = M(B, \bot)$.

\subsection{The two graphs are different in every closure-based sense}

\emph{Local clustering.} The local clustering coefficient of a vertex
$v$ of degree $d \ge 2$ is the fraction of pairs of neighbours of $v$
that are joined by an edge. In $A$, every vertex has two neighbours,
and those neighbours are joined (they are the other two vertices of
the same triangle), so the local clustering is $1$ at every vertex.
In $B$, every vertex has two neighbours, and those two neighbours are
at distance two apart along the cycle, not joined, so the local
clustering is $0$ at every vertex. Averaged:
\[
\Pi_{\clos}(A) = 1, \qquad \Pi_{\clos}(B) = 0.
\]

\emph{Triangle count.} A triangle in a graph is a set of three
vertices any two of which are joined. $A$ contains exactly two
triangles, namely $\{1, 2, 3\}$ and $\{4, 5, 6\}$. $B$ contains none.

\emph{Laplacian spectrum.} The Laplacian $L_\gph$ of a graph $\gph$
is the $|\vertices| \times |\vertices|$ matrix with entries
\[
(L_\gph)_{uv} =
\begin{cases}
\deg_\gph(v) & \text{if } u = v, \\
-1 & \text{if } \{u, v\} \in \edges, \\
0 & \text{otherwise.}
\end{cases}
\]
For a cycle $C_n$ the eigenvalues of $L_{C_n}$ are $2 - 2\cos(2\pi k
/ n)$ for $k \in \{0, 1, \dots, n-1\}$. Substituting gives
\[
\mathrm{spec}(L_A) = \{0, 0, 3, 3, 3, 3\},
\qquad
\mathrm{spec}(L_B) = \{0, 1, 1, 3, 3, 4\}.
\]
The multiplicities of the zero eigenvalue count the connected
components of the graph, and match ($A$ has two, $B$ has one). The
non-zero spectra differ.

\emph{Laplacian pseudoinverse.} The Moore-Penrose pseudoinverse
$L_\gph^\dagger$ is defined by inverting the non-zero eigenvalues of
$L_\gph$ and leaving the zero eigenspace alone. Since the non-zero
spectra of $L_A$ and $L_B$ differ, so do $L_A^\dagger$ and
$L_B^\dagger$.

\emph{Random walk mixing.} A simple random walk on $A$ starting at a
vertex $v$ stays in $v$'s triangle for all time; on $B$ it explores
the whole cycle. Any random-walk-based statistic that measures
communication between the two natural partitions (each triangle on
$A$, the two-colouring of $B$) separates the two graphs.

The pair $(A, B)$ therefore satisfies $A \sim_{\wl} B$ but differs
under every closure statistic and every spectral statistic on the
graph. The next chapter formalises what this means for MPNN
estimators.
